\begin{apartats}

\apartat[1,25]{1}
La funció
\[
f(x)=x^3+ax^2+3
\]
té un extrem relatiu en $x=2$. Troba $a$, digues de quin tipus d'extrem es tracta i estudia la
monotonia de $f$.

\begin{solucio}
$f'(x)=3x^2+2ax$. Si hi ha extrem en $x=2$, aleshores $f'(2)=12+4a=0$, d'on $\boxed{a=-3}$.\\
Amb $a=-3$: $f'(x)=3x^2-6x=3x(x-2)$, que s'anul·la en $x=0$ i $x=2$.\\
$f$ creix a $(-\infty,0)$, decreix a $(0,2)$ i creix a $(2,+\infty)$: en $x=0$ hi ha un
\textbf{màxim relatiu} i en $x=2$, un \textbf{mínim relatiu}.\\
Comprovació: $f''(x)=6x-6$ i $f''(2)=6>0$, que confirma el mínim.
\end{solucio}

\begin{nomesllarg}
\apartat{0,75}
\textbf{Inventa.} Escriu una funció que tingui un màxim relatiu en $x=-1$ i cap altre extrem.
Justifica-ho.

\begin{solucio}
Per exemple, $g(x)=-(x+1)^2$.\\
$g'(x)=-2(x+1)$, que només s'anul·la en $x=-1$, i hi canvia de positiva a negativa: hi ha un
màxim i no n'hi pot haver cap més, perquè la derivada no s'anul·la enlloc més.
\end{solucio}
\end{nomesllarg}

\apartat[1,25]{0,75}
Quants extrems relatius pot tenir, com a màxim, una funció polinòmica de grau $4$? Raona-ho.

\begin{solucio}
Els extrems relatius surten dels punts on s'anul·la la derivada. Si $f$ té grau $4$, aleshores
$f'$ té grau $3$ i, com a màxim, tres arrels reals.\\
Per tant, com a màxim hi pot haver \textbf{tres extrems relatius}. Per exemple,
$f(x)=x^4-2x^2$ en té tres, i $f(x)=x^4$ només un.
\end{solucio}

\end{apartats}
